\documentclass[a4paper,12pt]{article}
\usepackage[T2A]{fontenc}
\usepackage[utf8]{inputenc}
\usepackage[russian]{babel}
\usepackage{amsmath,amssymb,mathtools}
\usepackage{helvet}
\renewcommand{\familydefault}{\sfdefault}
\usepackage{icomma}
\usepackage[a4paper,top=18mm,bottom=18mm,left=21mm,right=21mm]{geometry}
\usepackage{microtype}
\usepackage{tikz}
\usetikzlibrary{calc,arrows.meta,positioning,shapes.geometric}
\usepackage{tabularx,booktabs}
\usepackage{enumitem}
\usepackage[hidelinks]{hyperref}
\usepackage{parskip}
\usepackage{fancyhdr}
\usepackage{setspace}
\usepackage{xcolor}

\definecolor{accent}{HTML}{187A78}
\definecolor{darkaccent}{HTML}{195052}
\definecolor{textgray}{HTML}{303D3E}
\definecolor{lightline}{HTML}{A9C7C3}
\color{textgray}
\setstretch{1.16}
\setlength{\parindent}{0pt}
\setlist[enumerate]{leftmargin=*,itemsep=.8em,topsep=.5em}
\setlist[itemize]{leftmargin=*,itemsep=.25em,topsep=.25em}
\setlength{\headheight}{14pt}
\pagestyle{fancy}
\fancyhf{}
\renewcommand{\headrulewidth}{0pt}
\renewcommand{\footrulewidth}{.3pt}
\fancyfoot[L]{\footnotesize Лёвин А.~А.\,·\,Анастасия Павлова}
\fancyfoot[R]{\footnotesize\thepage}

\newcommand{\checkitem}[1]{\noindent\(\square\)\enspace #1\par}
\newcommand{\keyterm}[1]{{\color{darkaccent}\textbf{#1}}}
\newcommand{\smallrule}{\par\vspace{.15em}{\color{lightline}\rule{\linewidth}{.35pt}}\vspace{.15em}}
\newcommand{\answer}[2]{#1 & #2 \\}

\tikzset{
  flowstart/.style={ellipse,draw=accent,thick,minimum width=36mm,minimum height=11mm,align=center,font=\small},
  flowstep/.style={rounded corners=3pt,draw=accent,thick,align=center,text width=109mm,minimum height=16mm,inner sep=6pt,font=\small},
  flowchoice/.style={diamond,aspect=2.1,draw=accent,thick,align=center,text width=48mm,inner sep=3.2pt,font=\small},
  flowbranch/.style={rounded corners=3pt,draw=accent,thick,align=center,text width=56mm,minimum height=21mm,inner sep=5pt,font=\small},
  flowarrow/.style={-{Stealth[length=2.8mm]},thick,color=darkaccent}
}

\begin{document}
\setlength{\abovedisplayskip}{7pt}
\setlength{\belowdisplayskip}{7pt}
\setlength{\abovedisplayshortskip}{4pt}
\setlength{\belowdisplayshortskip}{5pt}
\begin{titlepage}
\thispagestyle{empty}
\begin{center}
\vspace*{38mm}
{\Large\color{accent}\bfseries ЧЕК-ЛИСТ\par}
\vspace{12mm}
{\Huge\color{darkaccent}\bfseries Процентные изменения\par}
\vspace{6mm}
{\LARGE\color{darkaccent}\bfseries и текстовые задачи\par}
\vspace{11mm}
{\Large Повышение и снижение цены, проценты\par}
{\Large от новой базы, повторение смесей и систем\par}
\vspace{22mm}
{\large Математика\par}
\vspace{11mm}
{\large 15.10.26\par}
\vfill
{\normalsize Лёвин Артём Александрович\par}
{\normalsize эксклюзивно для Анастасии Павловой\par}
\vspace{21mm}
\begin{tikzpicture}[remember picture,overlay]
 \draw[accent!48,line width=1.1pt]
  ([xshift=22mm,yshift=19mm]current page.south west)
  .. controls ([xshift=78mm,yshift=28mm]current page.south west) and
             ([xshift=-65mm,yshift=8mm]current page.south east) ..
             ([xshift=-22mm,yshift=19mm]current page.south east);
\end{tikzpicture}
\end{center}
\end{titlepage}
\setcounter{page}{2}

\section*{1. Как найти процент от числа}
\checkitem{\keyterm{Один процент} составляет сотую часть величины. Чтобы найти \(p\%\) от числа \(x\), умножьте \(x\) на \(\frac{p}{100}\):}
\[
\boxed{\frac{p}{100}\,x=rx,\qquad r=\frac{p}{100}.}
\]
\checkitem{\(r\) --- доля изменения в десятичной записи. Например, \(11\%=0{,}11\), \(4\%=0{,}04\), \(17\%=0{,}17\). Чтобы вернуться к процентам, умножьте \(r\) на 100.}
\textbf{Пример.} \(20\%\) от \(100\) равно \(0{,}2\cdot100=20\). А \(20\%\) от неизвестного \(x\) --- это \(0{,}2x\).

\smallrule
\section*{2. Одна таблица для задач на повышение и снижение}
\checkitem{Обозначьте исходную цену через \(x\), а долю изменения через \(r=p/100\). Таблица показывает, какая часть от \(x\) добавляется или вычитается.}
\begin{center}
\renewcommand{\arraystretch}{1.32}
\begin{tabularx}{\linewidth}{@{}Xccc@{}}
\toprule
\textbf{Изменение} & \textbf{Было} & \textbf{Изменилось на} & \textbf{Стало}\\
\midrule
Повышение на \(p\%\) & \(x\) & \(+rx\) & \(x+rx=(1+r)x\)\\
Снижение на \(p\%\) & \(x\) & \(-rx\) & \(x-rx=(1-r)x\)\\
\bottomrule
\end{tabularx}
\end{center}
\checkitem{При \keyterm{повышении} коэффициент \(1+r\) больше 1. При \keyterm{снижении} коэффициент \(1-r\) меньше 1 (если \(0<r<1\)).}
\checkitem{Когда известна конечная цена \(S\), а требуется первоначальная, составьте уравнение \((1\pm r)x=S\).}

\subsection*{Пример 1. Цена выросла на 11\%}
После повышения товар стоит 1443~руб. Сколько он стоил раньше?
\[
\underbrace{x}_{\text{было}}+\underbrace{0{,}11x}_{\text{рост}}=
\underbrace{1443}_{\text{стало}},\qquad
1{,}11x=1443,\qquad \boxed{x=1300.}
\]
\subsection*{Пример 2. Цена снизилась на 13\%}
После снижения товар стоит 1218~руб. Тогда
\[
x-0{,}13x=1218,\qquad 0{,}87x=1218,\qquad \boxed{x=1400.}\]
\checkitem{\keyterm{Типичная ошибка:} при снижении на 13\% использовать коэффициент \(1{,}13\). Верный коэффициент --- \(0{,}87\).}

\section*{3. Как найти сам процент изменения}
\checkitem{Если известны значения \(S_0\) \emph{до} изменения и \(S_1\) \emph{после}, сначала найдите разность, затем разделите её \keyterm{на первоначальное значение}:}
\[
\boxed{r=\frac{|S_1-S_0|}{S_0},\qquad p=100r\%.}
\]
\textbf{Пример.} Цена снизилась с 2800 до 2324~руб. Тогда
\[
r=\frac{2800-2324}{2800}=\frac{476}{2800}=0{,}17,
\qquad \boxed{p=17\%.}\]
\checkitem{\keyterm{Контроль:} 476~руб. составляют 17\% от 2800~руб., поэтому результат разумен.}

\smallrule
\section*{4. «В несколько раз» и «на несколько процентов»}
\checkitem{Если величина \keyterm{уменьшилась в \(k\) раз}, она стала равна \(x/k\). При \(k>1\) доля снижения определяется из уравнения}
\[
(1-r)x=\frac{x}{k}\ \ (x>0),\qquad
\boxed{r=1-\frac{1}{k},\qquad p=\left(1-\frac{1}{k}\right)100\%.}
\]
\textbf{Пример.} Себестоимость снизилась в \(2{,}5\) раза:
\[
 r=1-\frac{1}{2{,}5}=1-0{,}4=0{,}6,
\qquad \boxed{p=60\%.}
\]
\checkitem{\keyterm{Важно:} «уменьшилась в \(2{,}5\) раза» означает \emph{деление} на \(2{,}5\); «уменьшилась на \(2{,}5\%\)» означает умножение на \(0{,}975\).}

\section*{5. Два последовательных процентных изменения}
\checkitem{После первого изменения появляется \keyterm{новая база}. Следующий процент берут от нового значения; проценты последовательных изменений нельзя просто складывать.}
\[
\boxed{S_1=S_0(1+r_1),\qquad S_2=S_1(1+r_2).}
\]
\textbf{Пример.} В квартале проживали 40\,000 человек. Число жителей выросло сначала на 8\%, затем ещё на 9\% \emph{от нового числа жителей}:
\begin{center}
\renewcommand{\arraystretch}{1.2}
\begin{tabularx}{\linewidth}{@{}Xrrr@{}}
\toprule
\textbf{Этап} & \textbf{Было} & \textbf{Прирост} & \textbf{Стало}\\
\midrule
Первый год, +8\% & 40\,000 & 3\,200 & 43\,200\\
Второй год, +9\% & 43\,200 & 3\,888 & 47\,088\\
\bottomrule
\end{tabularx}
\end{center}
\[
\boxed{40\,000\cdot1{,}08\cdot1{,}09=47\,088.}
\]

\section*{6. Повторение: вычитание уравнений системы}
\checkitem{Если части уравнений совпадают, \keyterm{вычтите одно уравнение из другого}: одинаковые слагаемые сократятся.}
\textbf{Пример.} Пусть \(A=0{,}4x+0{,}7y\), \(T=x+y+10\). Тогда система
\[
\begin{cases}A=0{,}41T,\\A+6=0{,}53T\end{cases}
\quad\Longrightarrow\quad 6=0{,}12T\quad\Longrightarrow\quad T=50.
\]
Значит, \(x+y=40\). Подстановка \(y=40-x\) в первое уравнение даёт \(x=25\), \(y=15\). \keyterm{Проверьте оба уравнения.}

\smallrule
\section*{7. Повторение: удаление воды и добавление раствора}
\checkitem{Основная величина --- \keyterm{масса растворённого вещества}: \(m_{\text{в-ва}}=w\,m_{\text{раствора}}\). При удалении \emph{чистой воды} масса вещества не меняется.}
\textbf{Пример.} Из 30~кг 20\%-го раствора удалили \(x\)~кг чистой воды и добавили \(x\)~кг 60\%-го раствора. Итоговая концентрация --- 35\%.
\begin{center}
\renewcommand{\arraystretch}{1.2}
\begin{tabularx}{\linewidth}{@{}Xccc@{}}
\toprule
\textbf{Этап} & \textbf{Масса, кг} & \textbf{Доля} & \textbf{Вещества, кг}\\
\midrule
Исходный раствор & \(30\) & \(0{,}2\) & \(6\)\\
Удалённая вода & \(-x\) & \(0\) & \(0\)\\
Добавленный раствор & \(+x\) & \(0{,}6\) & \(0{,}6x\)\\
Итоговая смесь & \(30\) & \(0{,}35\) & \(10{,}5\)\\
\bottomrule
\end{tabularx}
\end{center}
\[
6+0{,}6x=10{,}5\quad\Longrightarrow\quad \boxed{x=7{,}5\ \text{кг}.}
\]
\checkitem{Проверьте физическую возможность: удалить можно не больше воды, чем было в исходном растворе. В этом примере воды первоначально \(30-6=24\)~кг.}

\clearpage
\thispagestyle{plain}
\begin{center}
{\Large\color{darkaccent}\bfseries Алгоритм решения процентной задачи}\par
\vspace{3mm}
{\small От чтения условия к проверке результата}\par
\vspace{9mm}
\begin{tikzpicture}
\node[flowstart] (a) at (0,0) {Начало};
\node[flowstep] (b) at (0,-1.55) {Определите, что было и что стало.\\ Переведите \(p\%\) в \(r=p/100\).};
\node[flowchoice] (c) at (0,-4.2) {Изменений несколько?};
\node[flowbranch] (d) at (-4.65,-7.15) {Одна строка таблицы:\\ \(S_1=S_0(1\pm r)\).};
\node[flowbranch] (e) at (4.65,-7.15) {Несколько строк:\\ новая база после каждого этапа.};
\node[flowstep] (f) at (0,-9.9) {Составьте уравнение и найдите неизвестное.\\ При необходимости переведите долю в проценты.};
\node[flowstep] (g) at (0,-12.1) {Проверьте знак, базу для процента,\\ единицы и разумность ответа.};
\node[flowstart] (h) at (0,-13.75) {Готово};
\draw[flowarrow] (a) -- (b);
\draw[flowarrow] (b) -- (c);
\draw[flowarrow] (c.west) -| node[pos=.27,above,font=\small]{нет} (d.north);
\draw[flowarrow] (c.east) -| node[pos=.27,above,font=\small]{да} (e.north);
\draw[flowarrow] (d.south) |- (f.west);
\draw[flowarrow] (e.south) |- (f.east);
\draw[flowarrow] (f) -- (g);
\draw[flowarrow] (g) -- (h);
\end{tikzpicture}
\end{center}


\clearpage
\section*{Тренировочный блок · Путь А}
\textit{Решите 10 заданий. Записывайте уравнения или таблицы там, где они помогают. В ответах приведены только итоговые результаты.}
\vspace{2mm}
\begin{enumerate}[label=\textbf{\arabic*.},itemsep=1.05em]
\item Найдите \(18\%\) от 2500. Запишите также выражение для \(7\%\) от произвольного числа \(x\).
\item После повышения цены на 12\% товар стал стоить 1680~руб. Сколько он стоил до повышения?
\item После снижения цены на 15\% товар стал стоить 1530~руб. Какова его первоначальная цена?
\item С учётом скидки 4\% покупатель заплатил 1824~руб. Сколько стоила покупка без скидки?
\item Цена товара снизилась с 3500 до 2870~руб. На сколько процентов снизилась цена?
\item Себестоимость детали уменьшилась в 4 раза. На сколько процентов она снизилась?
\item В посёлке жили 15\,000 человек. За первый год число жителей увеличилось на 8\%, за второй --- ещё на 5\% по сравнению с предыдущим годом. Сколько жителей стало после второго года?
\item Число жителей сначала увеличилось на 10\%, а затем ещё на 20\% от нового количества. В итоге стало 26\,400 жителей. Сколько их было первоначально?
\item В сосуде 40~кг 15\%-го раствора кислоты. Из него удалили \(x\)~кг чистой воды, затем добавили \(x\)~кг 45\%-го раствора той же кислоты. Концентрация стала равна 24\%. Найдите \(x\).
\item Решите систему уравнений и найдите \(x\) и \(y\):
\[
\begin{cases}
0{,}3x+0{,}5y=0{,}34(x+y+10),\\
0{,}3x+0{,}5y+4=0{,}42(x+y+10).
\end{cases}
\]
\end{enumerate}

\clearpage
\section*{Ответы к тренировочному блоку}
\begin{center}
\renewcommand{\arraystretch}{1.5}
\begin{tabularx}{.92\linewidth}{@{}cX@{}}
\toprule
\textbf{№} & \textbf{Ответ}\\
\midrule
\answer{1}{\(450;\quad 0{,}07x\)}
\answer{2}{\(1500\)~руб.}
\answer{3}{\(1800\)~руб.}
\answer{4}{\(1900\)~руб.}
\answer{5}{\(18\%\)}
\answer{6}{\(75\%\)}
\answer{7}{\(17\,010\) человек}
\answer{8}{\(20\,000\) человек}
\answer{9}{\(8\)~кг}
\answer{10}{\(x=15,\quad y=25\)}
\bottomrule
\end{tabularx}
\end{center}
\vspace{9mm}
\textit{Короткая самопроверка.} При повышении конечная величина должна быть больше исходной, при снижении --- меньше. В многоэтапных задачах проверьте, что каждый новый процент рассчитан от величины, полученной на предыдущем шаге.

\end{document}
